\def\PUDhabilidades{ \vm	
	\textbf{CG x / HG yy } – ;
	
	\textbf{CG x / HG yy} – ;
	
	\textbf{CG x / HG yy} –.
}

\def\PUDatitudes{ \vm
	\textbf{Crítica e ética (CG1/CG7):} Postura crítica e ética na manipulação e processamento de dados ... \\
	
}


\def\PUDcodigo{ACE010}

\if\COMPILAR0
   \def\PUDcodacad{04507.12}
\else
   \def\PUDcodacad{04506.12}
\fi

\def\PUDdisciplina{Cálculo 3}

\def\PUDsemestre{S3}

\def\PUDhoras{80}

%NOVOS ---------------------------------------------------
\def\PUDhorasTeoria{80}
\def\PUDhorasPratica{0}

\def\PUDinterECA{---}

\def\PUDatuaisECA{---}

\def\PUDrecursos{Projetor multimídia; Quadro branco e pincel.}

%----------------------------------------------------------

\def\PUDcreditos{4}

\def\PUDprerequisito{ \hyperref[PUD:ACE006]{ACE006} }

\if\COMPILAR0
   \def\PUDpreacad{ \hyperref[PUD:ACE006]{Cálculo 2 (04507.7)} }
\else
   \def\PUDpreacad{ \hyperref[PUD:ACE006]{Cálculo 2 (04506.7)} }
\fi

\def\PUDobjetivos{Compreender o ramo do Cálculo 3 através dos seguintes objetivos específicos: Deduzir equações paramétricas e simétricas da reta no espaço, utilizando vetores no R3;  Identificar as principais superfícies quádricas, como parabolóides, elipsóides e hiperbolóides;  Definir função de várias variáveis identificando domínio, imagens e curvas de nível;  Resolver integral dupla e tripla através da interpretação geométrica;  Calcular integrais de linha e de superfície, aplicando o teorema de Green em regiões simples.}

\def\PUDementa{Estudo de vetores, retas e planos no R3.  Superfície quádricas.  Funções de várias variáveis.  Diferenciação.  Integrais múltiplas.  Campos vetoriais.}

\def\PUDconteudo{ &  Vetores no R3: Vetores;  Produto interno;  Produto vetorial;  A reta no R3;  O plano no R3.  Vetores;  Produto interno;  Produto vetorial;  A reta no R3;  O plano no R3. 
\\ 
 &  Superfície quádricas: Equação geral das quádricas;  Parabolóide;  Elipsóide;  Hiperbolóide. 
\\ 
 &  Funções de várias variáveis: Definição;  Domínio e imagem;  Curvas de níveis;  Derivadas Parciais;  Diferenciabilidade e o diferencial total;  A regra da cadeia;  Derivada direcional e gradiente;  Plano tangente e reta normal;  Derivadas parciais de ordem superior;  Máximo e Mínimo relativo e Multiplicadores de Lagrange. 
\\ 
 &  Integrais Múltiplas: A integral dupla;  Integral Iterada;  Integrais duplas em coordenadas polares;  Integral tripla;  Integral tripla em coordenadas cilíndricas e esféricas. 
\\ 
 &  Campos Vetoriais: Funções vetoriais;  Campos conservativos e integral de linha;  O teorema de Green;  O rotacional e o divergente;  O teorema de Stokes. }

\def\PUDmetodo{Aulas expositórias que podem ser teóricas e/ou práticas, onde as práticas no laboratório serão marcadas no decorrer da disciplina.}

\def\PUDavalia{A avaliação será feita com aplicação de provas teóricas e/ou práticas, além da possibilidade de inclusão de trabalhos e seminários no decorrer da disciplina.}

\def\PUDbibbasica{ & \href{www.biblioteca.ifce.edu.br/index.asp?codigo_sophia=23723}{Guidorizzi}, Hamilton L.  Um curso de cálculo.  Volume 3.  Rio de Janeiro, RJ : LTC, 2002.  362p.  ISBN: 9788521612575. 
\\
 &   \href{www.biblioteca.ifce.edu.br/index.asp?codigo_sophia=65473}{Boulos}, Paulo.  Calculo diferencial e integral.  2º ed.  São Paulo, SP: Pearson Education do Brasil, 2002.  v. 02.  349p.  ISBN: 9788534614580. 
\\
&   \href{www.biblioteca.ifce.edu.br/index.asp?codigo_sophia=61965}{Hoffmann}, Laurence D.  Calculo: um curso moderno e suas aplicações.  11º ed.  Rio de Janeiro, RJ: LTC, 2016.  661p.  ISBN: 9788521625315. }
%&   SWOKOSKI, Earl W.  Cálculo com geometria analítica.  Tradução de Alfredo Alves de Farias.  2.  ed.  São Paulo: Makron Books, 1997.  v. 2.  ISBN :8534603103. }

\def\PUDbibcomplementar{ &  \href{www.biblioteca.ifce.edu.br/index.asp?codigo_sophia=62611}{Hughes}-Hallett, Deborah.  Calculo Aplicado.  4º ed.  Rio de Janeiro, RJ: LTC, 2012.  483p.  ISBN: 9788521620518.  
\\ 
 &  \href{www.biblioteca.ifce.edu.br/index.asp?codigo_sophia=62607}{Anton}, Howard.  Cálculo.  10º ed.  Porto Alegre, RS: Bookman, 2014.  v. 2.  669p.  ISBN: 9788582602454.  
\\ 
 &  \href{www.biblioteca.ifce.edu.br/index.asp?codigo_sophia=24122}{Gonçalves}, Mirian Buss.  Cálculo B: funções de várias variáveis, integrais múltiplas, integrais curvilíneas e de superfície.  2º ed.  São Paulo, SP: Pearson Prentice Hall, 2007.  435p.  ISBN: 9788576051169.
\\
 &  \href{www.biblioteca.ifce.edu.br/index.asp?codigo_sophia=58165}{Fernandes}, Daniela Barude.  Cálculo Diferencial.  E-book.  Pearson.  132p.  ISBN: 9788543005423.

\\
 &  \href{www.biblioteca.ifce.edu.br/index.asp?codigo_sophia=55377}{Weir}, Maurice D.; Hass, Joel; Giordano, Frank R.  Cálculo: George B. Thomas.  E-book. v.2.  11º ed.  Pearson.  664p.  ISBN: 9788588639362. }

\def\PUDautor{Samuel Vieira Dias}

\def\PUDdata{2013-04-23}

\def\PUDrevisao{1}

\def\PUDrevisor{Samuel Vieira Dias}

\def\PUDdatarevisao{2017-05-09}

\PUDcorpo
